\section{\bi{Background and attribution}{背景與研究歸屬}}
\bi{A hard minimum spacing changes what a dense alternating chain can do. At zero spacing an adjacent positive and negative charge can cancel exactly. Once a positive floor is imposed, the residual interaction of the almost neutral pairs can select both the number of particles in each cluster and the positions of the individual particles. We ask for all minimizers at fixed period and particle count. Three different comparisons will matter: against a single trial cell, over the whole family of two packed triples, and over the entire feasible simplex.}{在稠密的交替鏈中，相鄰正負電荷可在零間距時完全抵消。給定正的硬下限後，近中性粒子對仍有相互作用，既可選出簇內的粒子數，也可固定個別粒子的位置。本文在固定週期及粒子數下求全部極小者。與單一試探構型比較、在整個緊雙三簇族取最小值，以及在全部可行間距單純形取最小值，是三個不同的命題。}

\bi{Bétermin, Knüpfer and Nolte introduced the alternating-chain energy and posed the equidistance conjecture for completely monotone profiles in the squared distance \cite{bkn}. The upstream manuscript gives a rigorous four-site Gaussian counterexample \cite{base}. A subsequent Lean proof compares the continuous four-site displacement family with the uniform cell for $0<q\le1/4$ \cite{prseven}. The dilute theorem \cite{dilute} proves equidistance for $\alpha(L/N)^2\ge4$. Its parameter is the squared dimensionless mean spacing. Here the period and particle count fix that mean while the independent floor can vary.}{Bétermin、Knüpfer 與 Nolte 定義交替鏈能量，並對距離平方的完全單調勢提出等距猜想 \cite{bkn}。上游論文已有嚴格的四粒子高斯反例 \cite{base}。其後的 Lean 證明在 $0<q\le1/4$ 比較了連續四粒子位移族與均勻構型 \cite{prseven}。稀薄區定理 \cite{dilute} 證明 $\alpha(L/N)^2\ge4$ 時的等距結晶，其參數為無因次平均間距的平方。本題的平均間距由週期及粒子數固定，硬下限則可獨立改變。}

\bi{At $q\le1/3$, the infimum manuscript \cite{infimum} established the collision dipole, semicircle stability, localization of each particle, the balanced odd split and the quadratic small-floor cost. Its trial state already had packed antipodal clusters. The four-site theorem recorded with Issue 5 \cite{four} proved exact packing at $s=4$, $0<h\le1/(8\pi)$ in our units. It used gap derivatives to obtain equality. Its larger-cell tests include centered triples and the five-contact trial. Circle rearrangement \cite{circle} and stability theory on compact groups \cite{riesz} provide the analytic background. Square-root displacement laws have precedents for bipartite trial lattices of a single species \cite{bipartite}.}{碰撞下確界論文 \cite{infimum} 已在 $q\le1/3$ 證明對徑偶極、半圓穩定性、逐粒子定位、均衡奇數分簇與小間距的二次能量代價，所採用的試探構型已是對徑緊雙簇。Issue 5 所附實驗室稿件的四粒子定理 \cite{four} 利用間距導數確定取等；換算本文單位，其範圍為 $s=4$、$0<h\le1/(8\pi)$。該研究的較大粒子數實驗包括居中三粒子簇與五接觸試探。圓周重排 \cite{circle} 和緊群穩定性理論 \cite{riesz} 提供解析背景；單一粒子種類的雙亞格子試探族已有平方根位移律的先例 \cite{bipartite}。}

\bi{The candidate contributions relative to these records are the quantified exact contact and antipodal locking for general even counts, the fourth-order coefficient, the global six-site phase strip including its two transition curves and widening interval, and a different global ground state at one outlying parameter pair. Attribution is relative to the fixed manuscripts and literature records preserved with the source index.}{相對於上述文獻，本項目可主張的候選增量為一般偶數粒子數的定量精確接觸與對徑鎖定、四次係數、六粒子的全域三相條帶及兩條相界與相寬，以及條帶外一點的另一種全域基態。原創性判斷限於研究索引凍結時本地保存的手稿與文獻記錄。}

\section{\bi{Model, normalization and equivalence}{模型、歸一化與等價關係}}
\bi{Fix $N=2m$, $L>0$ and $\alpha>0$. Write $x_j=X_j/L$, $x_{j+N}=x_j+1$, $\sigma_j=(-1)^j$, $g_j=x_{j+1}-x_j$. All $N$ gaps, including the last, satisfy $g_j\ge h=h_{\rm phys}/L$, $\sum g_j=1$. The mean gap is $1/N$. We use}{固定 $N=2m$、$L>0$ 及 $\alpha>0$。令 $x_j=X_j/L$、$x_{j+N}=x_j+1$、$\sigma_j=(-1)^j$、$g_j=x_{j+1}-x_j$。包括跨週期間隙在內，全部 $N$ 個間隙均滿足 $g_j\ge h=h_{\rm phys}/L$、$\sum g_j=1$。平均間距為 $1/N$。記}
\begin{equation}\label{eq:model}
 s=\alpha L^2,\quad q=\ee^{-\pi^2/s},\quad c=\sqrt{\pi/s},\quad
 \theta_j=2\pi x_j,\quad \eta=2\pi h,\quad
 G_s(d)=\sum_{p\in\mathbb Z}\ee^{-s(d+p)^2}.
\end{equation}
\bi{The ordered energy per particle, and its Fourier representation by Poisson summation, are}{對每粒子計的有向能量及由 Poisson 求和得到的 Fourier 表示為}
\begin{equation}\label{eq:energy}
 E_s=1-\frac1N\sum_{i,j}\sigma_i\sigma_j G_s(x_j-x_i)
 =1-\frac{2c}{N}\sum_{n\ge1}q^{n^2}|P_n|^2,\qquad
 P_n=\sum_j\sigma_j\ee^{in\theta_j}.
\end{equation}
\bi{Neutrality removes the constant Fourier mode. The leading 1 removes exactly the central self interaction per particle. Images with the same label in other cells remain. Total energy is $NE_s$. The unordered-pair score $W=\sum_{i<j}\sigma_i\sigma_j G_s(x_j-x_i)$ satisfies $E_s=1-G_s(0)-2W/N$. Derivative series converge absolutely and uniformly on compact parameter sets with $s>0$. For $0<h\le1/N$, the simplex modulo translation is compact. We write $M_N(h;s)$ for its attained minimum. At $h=0$, distinguish the open strictly ordered class from the compact collision closure. Rotations translate all sites. Even cyclic relabelings preserve species; an odd shift with the first sign reset reverses all charges. Reflections can preserve species; global sign reversal is a separate equivalence.}{電中性消去 Fourier 零模。常數 1 恰刪去每粒子的中央自身相互作用，其餘週期的同標號像仍計入。總能量為 $NE_s$；若 $W=\sum_{i<j}\sigma_i\sigma_j G_s(x_j-x_i)$ 對無序粒子對求和，則 $E_s=1-G_s(0)-2W/N$。各階導數級數於 $s>0$ 的緊參數集上絕對一致收斂。當 $0<h\le1/N$，模平移的可行單純形為緊集，最低值記作 $M_N(h;s)$。在 $h=0$ 時分別考慮嚴格有序類和容許碰撞的閉包。旋轉即全體平移。偶數步循環重標保留種類；奇數步重標並重置首電荷相當於整體換號。可另取保號反射，並明確區分是否容許整體換號。}

